% Arabisch nummerierte Abschnitte (Standard article)
\documentclass[11pt,a4paper]{article}

\usepackage[utf8]{inputenc}
\usepackage[T1]{fontenc}
\usepackage[ngerman]{babel}
\usepackage{amsmath,amssymb}
\usepackage{microtype}
\usepackage{algorithm}
\usepackage{algpseudocode}
\usepackage{hyperref}

\hypersetup{
  colorlinks=true,
  linkcolor=blue!60!black,
  urlcolor=blue!60!black,
  citecolor=blue!60!black,
}

\title{Entanglement-getemperte Zeitstempel-Spektroskopie:\\
  Quaternionisches EABC-Modell und der Bamberger Zeitwürfel}
\author{(Aus dem Python-Modul \texttt{Jitter Zeit.py})}
\date{\today}

\begin{document}
\maketitle

\begin{abstract}
  Wir beschreiben ein diskretes stochastisches Zeitmodell („Zeitwürfel“), in dem
  Verschränkung die Streuung der Zeitintervalle und eine effektive
  „Vierlingsbreite“ steuert. Zeitstempel werden modulo zwölf in vier aktive
  Zustände $E,A,B,C$ und einen Ruhezustand $X$ klassifiziert und als reelle
  Quaternionen in die Basis $\{1,i,j,k\}$ eingebettet. Zwischen aufeinanderfolgenden
  EABC-Ereignissen werden relative Drehungen („Schwung“), Übergangsmatriken und
  quaternionische Spektren mit komplexem Phasengewicht $\mathrm{e}^{-\mathrm{i}\Omega\Delta t}$
  ausgewertet. Der Text folgt der Implementierung und dient als begleitende
  Dokumentation; er erhebt keinen Anspruch auf physikalische Ableitung aus
  Erstprinzipien.
\end{abstract}

\section{Motivation und Überblick}

Das Modell kombiniert drei Bausteine:
\begin{enumerate}
  \item \textbf{Zufällige Zeitintervalle} $\Delta t_k$ um einen durch Temperatur $T$ und
    Verschränkungsparameter $S$ bestimmten Mittelwert, mit gaußschem Jitter.
  \item \textbf{EABC-Klassifikation} ganzer Zahlen $n_k \approx \alpha\, t_k$ modulo $12$
    mit Zuordnung zu Quaternion-Basisvektoren.
  \item \textbf{Spektralanalyse} der vier Kanäle $(q_k)_e,(q_k)_a,(q_k)_b,(q_k)_c$ als
    gewichtete Fourier-Summen über $\Delta t_k$ (bzw.\ aktive Zeiträume).
\end{enumerate}

\section{Quaternionische Darstellung der EABC-Zustände}

Sei $\mathbb{H}$ die Quaternionenalgebra mit Hamilton-Produkt,
$\mathrm{i}^2=\mathrm{j}^2=\mathrm{k}^2=\mathrm{ijk}=-1$.
Wir identifizieren die vier aktiven diskreten Zustände mit der Orthonormalbasis
in $\mathbb{R}^4 \cong \mathbb{H}$:
\begin{equation}
  E \leftrightarrow 1,\qquad
  A \leftrightarrow \mathrm{i},\qquad
  B \leftrightarrow \mathrm{j},\qquad
  C \leftrightarrow \mathrm{k}.
\end{equation}
Der Zustand $X$ (nicht in $\{1,5,7,11\}$ modulo $12$, siehe unten) wird als
Nullquaternion abgebildet.

\section{Klassifikation modulo $12$}

Für $n \in \mathbb{Z}$ sei $r = n \bmod 12$.
\begin{equation}
  \mathrm{label}(n) =
  \begin{cases}
    E, & r = 1,\\
    A, & r = 5,\\
    B, & r = 7,\\
    C, & r = 11,\\
    X, & \text{sonst.}
  \end{cases}
\end{equation}
In der Simulation wird $n_k = \mathrm{round}(\alpha\, t_k)$ gesetzt; $\alpha$ ist ein
freier Skalenparameter (\texttt{alpha\_scale}).

\section{Verschränkungsabhängiger Jitter und effektive Breite}

\subsection{Jitter-Standardabweichung}

Mit Verschränkungsparameter $S>0$, Boden $\sigma_{\min}$, Amplitude $\sigma_0$ und
Dämpfung $\gamma>0$:
\begin{equation}
  \sigma_t(S) = \sigma_{\min} + \sigma_0\,\mathrm{e}^{-\gamma S}.
\end{equation}

\subsection{Effektive Vierlingsbreite $\Delta_{\mathrm{eff}}$}

Glatte Variante:
\begin{equation}
  \Delta_{\mathrm{eff}}(T,S)
  = 8 + 2\bigl(1-\mathrm{e}^{-T/(2\pi)}\bigr)\bigl(1-\mathrm{e}^{-\lambda S}\bigr),
\end{equation}
$\lambda>0$. Grenzfälle: für $T\to 0$ gilt $\Delta_{\mathrm{eff}}\to 8$; für große $T,S$
nähert sich $\Delta_{\mathrm{eff}}$ dem Wert $10$.

Alternative („schärfer“) Temperaturabhängigkeit mit $\beta>1$:
\begin{equation}
  \Delta_{\mathrm{eff}}^{(\beta)}(T,S)
  = 8 + 2\bigl(1-\mathrm{e}^{-\beta T/(2\pi)}\bigr)\bigl(1-\mathrm{e}^{-\lambda S}\bigr).
\end{equation}

\section{Zeitwürfel-Simulation}

Gegeben $\omega_0>0$ und $\Delta_{\mathrm{eff}}$ wie oben setzt die Implementierung
\begin{equation}
  \Delta t^{(0)} = \frac{\Delta_{\mathrm{eff}}}{\omega_0}.
\end{equation}
Für $k=0,\ldots,N-1$:
\begin{equation}
  \varepsilon_k \sim \mathcal{N}(0,\sigma_t^2),\qquad
  \Delta t_k = \max\bigl\{\Delta t^{(0)}+\varepsilon_k,\;\Delta t^{(0)}\bigr\},
\end{equation}
\begin{equation}
  t_{k+1} = t_k + \Delta t_k,\quad t_0=0.
\end{equation}
Zu jedem Schritt werden $n_k$, $\mathrm{label}(n_k)$ und die Quaternion-Komponenten
von $q_k$ gespeichert.

\begin{algorithm}[htbp]
  \caption{Ablauf der Zeitwürfel-Simulation und nachgeschaltete Auswertung
    ($\Delta_{\mathrm{eff}}$ hier mit Temperaturfaktor $\beta$ wie in \texttt{delta\_eff\_sharp};
    in der Implementierung wird $\Delta t_k$ zusätzlich nach unten durch $\Delta_{\mathrm{eff}}/\omega_0$
    begrenzt, falls $\varepsilon_k$ zu negativ ausfällt).
    Die kompakte Schreibweise $\mathcal{P}_{\mathbb{H}}$ entspricht im Code der Summe der kanalweisen
    Leistungen $|Q_X|^2$, $X\in\{e,a,b,c\}$, vgl.\ Abschnitt~\ref{sec:spektrum}.}
  \label{alg:zeitwuerfel}
  \begin{algorithmic}[1]
    \State Wähle $N,T,S,\omega_0,\sigma_{\min},\sigma_0,\gamma,\lambda,\beta$.
    \State Berechne den entanglement-abhängigen Jitter
    \[
      \sigma_t(S)=\sigma_{\min}+\sigma_0\,\mathrm{e}^{-\gamma S}.
    \]
    \State Berechne die effektive Breite
    \[
      \Delta_{\mathrm{eff}}(T,S)
      =
      8+
      2\bigl(1-\mathrm{e}^{-\beta T/(2\pi)}\bigr)\bigl(1-\mathrm{e}^{-\lambda S}\bigr).
    \]
    \For{$k=1,\ldots,N$}
      \State Ziehe $\varepsilon_k\sim\mathcal{N}(0,\sigma_t^2)$.
      \State Setze
      \[
        \Delta t_k=\Delta_{\mathrm{eff}}/\omega_0+\varepsilon_k.
      \]
      \State Setze
      \[
        t_k=\sum_{j\leq k}\Delta t_j.
      \]
      \State Runde $n_k=\lfloor \alpha\, t_k\rceil$ und klassifiziere $n_k\bmod 12$.
      \State Weise
      \[
        E\mapsto 1,\quad A\mapsto \mathrm{i},\quad B\mapsto \mathrm{j},\quad C\mapsto \mathrm{k}
      \]
      zu (Einbettung von $q_k\in\mathbb{H}$).
    \EndFor
    \State Filtere die aktive EABC-Folge.
    \State Berechne den Schwung $u_k=\overline{q_k}\,q_{k+1}$.
    \State Berechne das Spektrum
    \[
      \mathcal{P}_{\mathbb H}(\Omega)
      =
      \left|
      \sum_k q_k\,\mathrm{e}^{-\mathrm{i}\Omega\Delta t_k}
      \right|^2.
    \]
  \end{algorithmic}
\end{algorithm}

Für zwei Quaternion-Zustände $q_k,q_{k+1}$ der \emph{vollen} Simulation definiert der Code
den relativen Schwung
\begin{equation}
  u_k := \overline{q}_k\, q_{k+1} \in \mathbb{H},
\end{equation}
mit Konjugation $\overline{q}$ und Hamilton-Produkt. Die Komponenten $(u_k)_e,(u_k)_a,(u_k)_b,(u_k)_c$
und $\|u_k\|$ werden protokolliert.

\section{Aktive EABC-Folge und Übergänge}

Alle Zeilen mit $\mathrm{label}\in\{E,A,B,C\}$ bilden die \textbf{aktive Folge}.
Zwischen aufeinanderfolgenden aktiven Ereignissen werden Zeiten $t_k$ und
\begin{equation}
  \Delta t_{\mathrm{active},k} := t_k - t_{k-1}
\end{equation}
(wie in der Implementierung: erste Zeile \texttt{NaN}) verwendet.

Nur zwischen aufeinanderfolgenden \emph{aktiven} Quaternionen wird erneut der Schwung
\begin{equation}
  u = \overline{q}^{(\mathrm{akt})}_{k}\, q^{(\mathrm{akt})}_{k+1}
\end{equation}
gebildet.

\subsection{Schwung-Label}

Sei $u = u_e + u_a\mathrm{i}+ u_b\mathrm{j}+ u_c\mathrm{k}$.
Das Label wählt den Index $\mathrm{E},\mathrm{A},\mathrm{B},\mathrm{C}$ mit maximalem
$|\cdot|$ unter den reellen Koeffizienten und ein Vorzeichen:
\begin{equation}
  \texttt{schwung\_label} \in \{0,\;+\mathrm{E},\;-\mathrm{A},\;\ldots\}.
\end{equation}

\subsection{Übergangsmatrix und Idealrotation}

Die Matrix der Paare $(\mathrm{label}_{\mathrm{from}},\mathrm{label}_{\mathrm{to}})$
wird auf feste Zeilen und Spalten $E,A,B,C$ expandiert (fehlende Einträge gleich null).

Als \textbf{Idealrotation} zählt die Implementierung die Übergänge
\begin{equation}
  A\!\to\! E,\quad E\!\to\! C,\quad C\!\to\! B,\quad B\!\to\! A;
\end{equation}
ihr relativer Anteil an allen aktiven Übergängen wird ausgegeben (\texttt{Idealrotations-Anteil}).

\section{Quaternionisches Spektrum}
\label{sec:spektrum}

Sei $\Omega\in\mathbb{R}$ eine Kreisfrequenz und $\Delta t_k$ die gewählte Intervallfolge
(z.\,B.\ Simulationsintervall oder $\Delta t_{\mathrm{active}}$).
Mit Phasen $\phi_k(\Omega)=\mathrm{e}^{-\mathrm{i}\Omega \Delta t_k}$ definieren wir
komponentenweise
\begin{equation}
  Q_X(\Omega) := \sum_k (q_k)_X\,\phi_k(\Omega),\quad X\in\{e,a,b,c\},
\end{equation}
und die Leistungen $P_X(\Omega)=|Q_X(\Omega)|^2$ sowie
\begin{equation}
  P_{\mathrm{ges}}(\Omega) = P_E+P_A+P_B+P_C.
\end{equation}
Das entspricht einer komponentenweisen komplexen Fourier-Summe mit Gewichtung durch die
Zeitintervalle (nicht durch gleichmäßige Abtastung in $t$).

\section{Zentriertes Spektrum}

Um den Gleichanteil der Kanäle zu unterdrücken, ersetzt die Implementierung
\begin{equation}
  \tilde{q}_k := q_k - \overline{q},\qquad
  \overline{q}_X := \frac{1}{N'}\sum_k (q_k)_X,
\end{equation}
($N'$ Anzahl der verwendeten Zeilen nach Filter) und wendet dieselbe Spektralformel auf
$\tilde{q}$ an. Optional wird nur die aktive Teilfolge verwendet; für $\Delta t$ kann
man $\Delta t_k$ der Simulation oder $\Delta t_{\mathrm{active}}$ wählen (\texttt{dt\_column}).

\section{Rigidity-Proxy und Vergleichsfolgen}

Die Datei enthält weiterhin einen einfachen Autokorrelations-Proxy für normierte
Intervallfluktuationen sowie Hilfsfunktionen zum Vergleich zweier Intervallstatistiken
(vgl.\ Nullstellenabstände, Primlücken usw.). Details sind der Implementierung zu entnehmen.

\section{Implementierung}

Die beschriebenen Objekte entsprechen den Funktionen und Klassen in
\texttt{Jitter Zeit.py}, insbesondere:
\texttt{simulate\_zeitwuerfel},
\texttt{quaternionic\_spectrum},
\texttt{active\_eabc\_sequence},
\texttt{active\_transition\_schwung},
\texttt{transition\_matrix},
\texttt{centered\_quaternionic\_spectrum},
\texttt{print\_transition\_report}.

\section*{Hinweis}

Das Modell ist als exploratives und reproduzierbares Rechenexperiment gedacht.
Physikalische Interpretation von $T,S,\Omega$ und der EABC-Zuordnung bleibt dem Nutzer
überlassen.

\end{document}
